\section{n4. Explicit actual-primary threefold obstruction: generic sections on exceptional surfaces do not imply all-curve-nef divisorial parts.}
\textit{Supports: G331, n1.}

Precise THREE-dimensional obstruction to the surface competitor method; this is NOT a continuous-envelope counterexample. Support: ENTIRE qualified G331 and own n1 for the finite section/ACTUAL primary presentation argument. Let R=k[x,y,z]\_(x,y,z) for ANY field k, m=(x,y,z), and I=(z)+m\textasciicircum{}2, a proper primary ideal. Blow up the origin, with exceptional E\_0=P\textasciicircum{}2\_k. Blow up the two distinct k-rational points P=[1:0:0], Q=[0:1:0] of E\_0. Their disjoint point ideals can be blown up at once; equivalently do the two point blowups successively. All charts are regular (ordinary point blowups in regular coordinates). Let Y be the result, E\_1 the strict transform of E\_0, and E\_2,E\_3 the new exceptional planes. Then E\_1=Bl\_(P,Q)(P\textasciicircum{}2), while E\_2,E\_3 are disjoint copies of P\textasciicircum{}2.

This explicit Y is an ACTUAL primary model. Its second-stage negative tautological bundle O(-E\_2-E\_3) is relatively ample for Y-\textgreater{}Bl\_m(R); this is the blowup of the disjoint union of the two point ideals, invertible/unit on the complement. The first negative tautological bundle O(-E\_0) is ample. The COMPLETE immediately read Stacks29.38.7--8 gives, for an integer a sufficiently large, O(-E\_2-E\_3-a pullback(E\_0)) ample over R. Here pullback(E\_0)=E\_1+E\_2+E\_3. Thus E=aE\_1+(a+1)E\_2+(a+1)E\_3 is ACTUAL effective Cartier, full exceptional support, and O(-E) relatively ample. Stacks37.53.6 gives H\textasciicircum{}0(Y,O)=R. The finite affine-section localization/projective immersion/Rees argument proved in n1 applies verbatim in this dimension to THIS explicit Y: it requires normal target, ample negative exceptional bundle and projectivity, not surface resolution. A finite set of sections of O(-nE) gives a proper primary J and identifies Y with Bl\_J(R). This is an explicit construction and no cofinality assertion for arbitrary threefold modifications.

The divisorial part of I on Y is F\_I=E\_1+2E\_2+2E\_3. The order at E\_0 is one, realized by z. In the x-chart of the first blowup, with v=y/x and w=z/x, I O=x(w,x). At P the regular local parameters are x,v,w. Thus the residual ideal (w,x) has point order one, so the second exceptional order is two. The y-chart gives the identical calculation at Q. The strict E\_1 order remains one. All ideal generators are global sections of O(-F\_I) by normal height-one intersection, and at the generic point of EACH E\_i at least one realizes the minimum and restricts to a nonzero section, exactly as in the surface competitor proof.

Nevertheless A=-F\_I fails the mandatory contracted-curve positivity test. On E\_1 let H be the pullback of a line in P\textasciicircum{}2 and e\_P,e\_Q its two exceptional curves. Restricting divisors gives E\_1|\_(E\_1)=-H-e\_P-e\_Q, E\_2|\_(E\_1)=e\_P and E\_3|\_(E\_1)=e\_Q: the strict first divisor is pullback(E\_0)-E\_2-E\_3, and O(E\_0)|\_(E\_0)=O\_P2(-1). Hence
O(-F\_I)|\_(E\_1) has class H-e\_P-e\_Q.
Let L be the strict transform of the line z=0 through P and Q. It is a proper integral contracted curve with class H-e\_P-e\_Q. The usual point-blowup chart calculation gives H\textasciicircum{}2=1, H.e\_P=H.e\_Q=e\_P.e\_Q=0 and e\_P\textasciicircum{}2=e\_Q\textasciicircum{}2=-1; these are local normal-bundle degrees, with residue degrees one. Therefore deg\_L O(-F\_I)=1-1-1=-1. Equivalently the nonzero restricted generator z determines an effective divisor whose fixed component is the negative curve L. A nonzero section on the exceptional SURFACE does not imply nonnegative degree on that curve.

On E\_2 and E\_3 the restrictions of -F\_I are O(1), since E\_1 restricts to a hyperplane and each new exceptional normal bundle is O(-1). Nonnegativity there cannot repair the degree -1 on L in E\_1. Thus positivity was checked on an ACTUAL primary model, and a SPECIFIC actual contracted integral curve violates it; no purely numerical abstract matrix was substituted.

The every-support inequality -v(I)\textless{}=f\_(-F\_I) still holds: the ideal-sheaf inclusion I O\_Y subset O(-F\_I) pulls to the valuation ring of Frac(R/p) for each support p, with generic centre outside the exceptional fibre, and regular ratios give the inequality even for killed generators. What fails is the subsequent inference that f\_(-F\_I) is an all-curve-nef geometric-core competitor on Y. In contrast -v(I) has a negative effective relatively ample Cartier presentation on ACTUAL Bl\_I(R). The example therefore pinpoints a false step in the surface strategy while preserving the positivity of genuine primary profiles and existing smooth characteristic-zero CE. It gives no failure of CE in any characteristic and no theorem about arbitrary higher-dimensional model cofinality.

The intersection numbers are proved by the explicit blowup geometry, not by a program. Own finite\_diagnostics.py only evaluates the recorded integer degree expression and divisor coefficients; it is a diagnostic. Artifacts n4.md, n4.tex and reports/finite-diagnostics.json; rerun python finite\_diagnostics.py; python assemble.py; python verify.py.

\subsection*{Expanded threefold degree obstruction}
For $I=(z)+\mathfrak m^2$ on the explicit regular threefold model,
\[
F_I=E_1+2E_2+2E_3,\quad
(-F_I)|_{E_1}=H-e_P-e_Q,\quad L\sim H-e_P-e_Q,
\]
\[
\deg_L\mathcal O_Y(-F_I)=1-1-1=-1.
\]
The ample actual primary presentation uses
$E=aE_1+(a+1)E_2+(a+1)E_3$ for a sufficiently large integer $a$,
proved by relative composition ampleness and the finite section/Rees
argument. The nonzero generic sections on the $E_i$ therefore fail the
all-contracted-curve test on an actual primary model, not an abstract
unrealized divisor.

